\section{Local trials and the ordinary sweep}\label{sec:local-sweeps}

The increment estimates of Section~\ref{sec:sweep} concern a fixed
prestate and an arbitrary set of possible next ports.  We now turn them
into a complete ordinary sweep.  The obstacle is the volume of the torus:
it is much larger than the inverse damping scale.  We therefore expose
the next digits, and finally the remaining suffixes, in tiles and monitor each trial in a bounded
neighborhood of its tile.  Resolvent decay will show that these finitely
many tests control the global matrices as well.

Throughout Section~\ref{sec:local-sweeps}, the dimension \(d\ge3\) is fixed.  Constants
may depend on \(d\), but are independent of the torus volume and the
stage.  Terminal sites remain fixed.  We assume that
their collapsed distance, defined below, satisfies
\begin{equation}\label{eq:local-metric-input}
 d_T(i,j)\ge\tfrac12\dist(i,j)-t^{-1.03}.
\end{equation}
Here \(\dist\) is nearest-neighbor distance on the torus.  The estimate is
initially immediate, and Lemma~\ref{lem:clearing-locality} will preserve it
when terminal sites are added.

The analytic estimates use the following conditional interface.  If
\(\mathcal F\) is the full retained history before the next slot is
chosen and \(x\) is the selected port, its real scalar update satisfies
\[
 \EE(\xi_0\mid\mathcal F,x)=0,\qquad
 \EE(\xi_0^2\mid\mathcal F,x)=t,\qquad
 |\xi_0|\le C\sqrt t.
\]
The next exposed slot of a trial is conditionally uniform among its
unexposed physical sites.  No symmetry or two-point support is used for
the scalar.  For the final sweep,
\(t=t_m=2\eta\) and \(w=1/2\): an untouched true suffix \(\xi\)
has variance \(\eta=wt\), mean zero, and
\(|\xi|\le\sqrt{3wt}\), so \(\xi_0=\xi/\sqrt w\) has exactly
these bounds.  Applying this interface to a true suffix requires that
the suffix is still unqueried at its retained pre-query history.
Lemma~\ref{lem:conditional-freshness} supplies the conditional laws for
unqueried coordinates and independently supplied order variables.  The
trial rule below verifies that ordinary sweeps query such inputs.

\subsection{Decay and the inverse of \texorpdfstring{$I-K$}{I-K}}

Assign cost zero to each nearest-neighbor edge meeting a terminal site,
and cost one to every other nearest-neighbor edge.  The minimum path
cost is the pseudometric \(d_T\).  Ports at the same physical site have
distance zero in both metrics.
The following proof uses the exponential-conjugation method of Combes
and Thomas~\cite{CombesThomas1973}, applied to this collapsed metric.

\begin{lemma}[Decay in the collapsed distance]\label{lem:locality}
Suppose the true terminal principal matrix is invertible.  Let \(g\)
be the active inverse with imaginary diagonal damping between \(ct\)
and \(Ct\).  There are constants \(a,C>0\), depending only on these
bounds and \(d\), such that
\begin{equation}\label{eq:local-ct}
 |h_{ij}|\le Ct^{-1/2}\exp\{-atd_T(i,j)\}.
\end{equation}
This also holds in the port representation, with changed constants.
If \eqref{eq:local-metric-input} holds and \(u=t^3\), then
\begin{equation}\label{eq:local-physical-decay}
 |h_{ij}|\le Ct^{-1/2}\exp\{-cu\dist(i,j)\}.
\end{equation}
All constants are independent of the torus size.
\end{lemma}

\begin{proof}
Write the physical active inverse as \(g=M^{-1}\), where the real part
of \(M\) is the terminal Schur matrix.  Its dissipative part gives
\(\|M^{-1}\|\le C/t\).  Every off-diagonal entry created by eliminating
a connected terminal component joins sites at collapsed distance zero.
Consequently conjugation by
\(Q(i)=\exp\{atd_T(z,i)\}\) leaves all these entries unchanged.
Only direct active-active hopping entries change.  There are at most
\(2d\) in each row and column, and their relative change is \(O(at)\).
Thus \(\|QMQ^{-1}-M\|\le C_d at\).  For a sufficiently small fixed \(a\),
a Neumann series gives
\(\|QM^{-1}Q^{-1}\|\le C/t\).  Set \(z=j\) and multiply by
\(\sqrt t\) to obtain \eqref{eq:local-ct}.  The bounded port weights
preserve this estimate.

For completeness, the additive error in
\eqref{eq:local-metric-input} cannot simply be absorbed into a
polynomial prefactor after exponentiation.  Split instead at ordinary
distance \(t^{-2}\).  Below this distance, the crude bound
\(Ct^{-1/2}\) suffices, since \(u t^{-2}=t\).  Above it,
\(d_T(i,j)\ge\dist(i,j)/4\) for small \(t\), and
\eqref{eq:local-ct} is stronger than
\eqref{eq:local-physical-decay}.
\end{proof}

\begin{lemma}[Absolute row bound]\label{lem:inverse-row}
For an active inverse or its port representation as in
Lemma~\ref{lem:locality}, assume \eqref{eq:local-metric-input}, bounded positive \(p\),
\(\|K\|\le\rho_0<1\), and
\[
 Fp=q,\qquad q_i=p_i+O(t^\beta),\qquad \beta>0.
\]
Then, for sufficiently small \(t\),
\begin{equation}\label{eq:local-p-row}
 \|P\|_{\infty\to\infty}\le C\log(1/t),
 \qquad P=(I-K)^{-1}.
\end{equation}
Moreover \(|P_{ij}|\le C\exp\{-cu\dist(i,j)\}\), where \(u=t^3\)
and the constant \(c>0\) may be decreased.
\end{lemma}

\begin{proof}
The Ward identity and bounded positive \(p\) give
\(\sup_i\sum_jF_{ij}\le C\).  Since \(F\) is symmetric, the column
sums obey the same bound.  Tilt \(K\) by
\(Q(i)=e^{u\dist(z,i)}\).  Each row and column sum of
\(QKQ^{-1}-K\) is bounded by
\[
 \sup_i\sum_jF_{ij}\big(e^{u\dist(i,j)}-1\big).
\]
The contribution from distances at most \(t^{-2}\) is \(O(t)\).
For greater distances, the stronger estimate
\(|h_{ij}|\le Ct^{-1/2}e^{-ct\dist(i,j)}\) from the proof of
Lemma~\ref{lem:locality} makes the sum exponentially small.  The
row-column norm bound therefore gives
\(\|QKQ^{-1}\|\le\rho<1\), uniformly over \(z\).
Hence
\[
 |(K^k)_{ij}|\le\rho^k e^{-u\dist(i,j)}.
\]
Summing in \(k\) proves the entrywise assertion.

To improve the absolute row sum, use a different estimate for short
powers.  Entrywise \(|K^k|\le F^k\), and
\(Fp\le(1+Ct^\beta)p\), so
\[
 \sum_j|(K^k)_{ij}|\le C(1+Ct^\beta)^k.
\]
These sums are uniformly bounded for
\(k\le k_0=\lceil C_1\log(1/t)\rceil\).  For the other powers, use
\[
 \sum_j e^{-u\dist(i,j)}\le C_d u^{-d}=C_d t^{-3d},
\]
which remains valid with two ports per site.  Indeed, choosing minimal
torus displacement representatives embeds each distance shell into
\(\mathbb Z^d\), where the shell of radius \(r\) has at most
\(C_d(1+r)^{d-1}\) sites.  Summing the geometric shell series gives
the displayed bound uniformly in the torus size.  Choose \(C_1=C_1(d)\)
so that \(t^{-3d}\rho^{k_0}\le1\).  The first \(k_0\) powers then cost \(Ck_0\),
and all remaining powers together cost \(C\).
\end{proof}

The transformed-lag cap in Section~\ref{sec:sweep} supplies precisely the hypothesis
\(q=p+O(t^\beta)\) of Lemma~\ref{lem:inverse-row}.  Lemma~\ref{lem:inverse-row} therefore
completes the inverse-row input used in the lag increment estimates.
We shall also need decay for a single prospective update, even if that
update violates one of the finer caps.

\begin{lemma}[A prospective ordinary update]\label{lem:local-prospective}
Suppose the sweep starts by duplicating a physical state satisfying the
complete stage-boundary bounds of Section~\ref{sec:state}, including
self-consistency, the broad bounds and comparators \eqref{eq:comparators},
the boundary cap \(\|K^{\rm phys}\|\le1/3\), the fourth-power row cap
\(H_*=t^{1+\delta}\), and an active physical-site count of at least
\(N/2\), as in Lemma~\ref{lem:duplication}.  In particular its initial
data include \(D_0\asymp W\) and the sharp duplicated comparison
\eqref{eq:sweep-sharp-duplicate}.  Assume the broad bounds on the frozen
\(r,p\), the fixed metric bound \eqref{eq:local-metric-input}, and both
current prefix conditions \eqref{eq:sweep-global-caps} and
\eqref{eq:sweep-row-lag-caps}, where \(\beta=.42\delta\).  For small
\(t\), these hypotheses give a fixed margin \(\|K\|\le\rho_0<1\),
bounded positive \(p\), and \(|g_{ii}-r_i|=O(t^\beta)\).
Perform one bit-port update with the
conditional moment and size bounds above.  Its poststate
exists and, for small \(t\), satisfies
\[
 \|K'\|\le\rho_1<1,\qquad
 q'_i=p_i+O(t^\beta),\qquad
 \sup_i\sum_j|h'_{ij}|^2\le C.
\]
It satisfies Lemmas~\ref{lem:locality} and \ref{lem:inverse-row}
without any assumption on the post-update fine caps.  The lag and row
proxy increments obey
\begin{equation}\label{eq:local-scalar-tails}
 |\Delta\widetilde U_i|+|\Delta R_i|
 \le t^{-C}e^{-cu\dist(i,x)},
\end{equation}
where \(R_i\) denotes either normalized fourth-power row proxy.
The entries of \(J_x\), and the edge weights of \(L_x\), obey
\begin{equation}\label{eq:local-matrix-tails}
 t^{-C}\exp\{-cu[\dist(i,x)+\dist(j,x)]\}.
\end{equation}
\end{lemma}

\begin{proof}
Reserve ports keep the physical damping at least \(ct\), so the
updated inverse exists.  By the single-update formula,
\[
 |g'_{ii}-g_{ii}|\le C|h_{ix}|^2/\sqrt t
 \le C\sqrt{H_*/t}=Ct^{\delta/2}.
\]
Since \(\delta/2>\beta\), this proves the new diagonal estimate.
The formula \(h'_{ij}=h_{ij}-\gamma h_{ix}h_{xj}\), with
\(|\gamma|\le C\), also gives
\[
 \sum_j|h'_{ij}|^2
 \le 2\sum_j|h_{ij}|^2
       +C|h_{ix}|^2\sum_j|h_{xj}|^2\le C.
\]
The inherited stage-start data and the current cap conditions are the
prestate hypotheses of Proposition~\ref{prop:sweep-increments};
Lemma~\ref{lem:inverse-row} supplies its inverse-row input.
That proposition gives \(\|J_x\|\le Ct^{\delta/4}\).  Write
\(I-\widehat K=\Schur_{\{x\}}(I-K)\).  The Schur contraction proved
after \eqref{eq:sweep-global-caps} gives
\(\|\widehat K\|\le\|K\|\le\rho_0\), while the exact next-step identity is
\[
 I-K'=\Schur_{\{x\}}(I-K)-J_x=I-\widehat K-J_x.
\]
Thus \(\|K'\|\le\rho_0+Ct^{\delta/4}\le\rho_1<1\) for small \(t\).
The terminal set has not changed, so its metric bound
has not changed either.  Lemmas~\ref{lem:locality} and
\ref{lem:inverse-row} now apply.

The explicit formulas for \(J_x,L_x\) contain a factor joining each
surviving endpoint to \(x\).  Lemma~\ref{lem:locality} therefore gives
\eqref{eq:local-matrix-tails}.  In the exact lag formula
\(\Delta\widetilde U=P'Z^x/\sqrt t\), the vector \(Z^x\) has the
same single-center decay by its definition and the prestate bound on
\(U\).  The entrywise bound on \(P'\) preserves that decay under
convolution.  The fourth-power proxy formulas are finite sums of the
same factors.  Each convolution loses at most a power of \(u^{-1}\),
because \(\sum_j e^{-cu\dist(i,j)}\le C_d u^{-d}\); decreasing the
exponent constant absorbs the convolution.  This proves
\eqref{eq:local-scalar-tails} with a polynomial prefactor whose degree
may depend on \(d\).  Crucially,
none of these deductions used a post-update fine cap.
\end{proof}

\subsection{A finite-dimensional test for the energy increment}

At ordinary stage \(n\), partition the torus into cubes of side
\(L_{\mathrm{tile}}=4^{6n}\asymp t^{-6}\).  Around each tile choose concentric cubes
\(Q,Q'\) of side \(3L_{\mathrm{tile}},5L_{\mathrm{tile}}\).  They fit in the torus, and each family
has uniformly bounded overlap.  Let \(\chi\) equal one on \(Q\),
vanish at the boundary of \(Q'\), and have nearest-neighbor slopes at
most \(C/L_{\mathrm{tile}}\).

Write \(\mathcal P_0\) for the ports present at the start of the
duplicated sweep.  The fixed real local test space \(\mathcal H_T\)
consists of pairs \((\alpha,v)\), where \(\alpha\) is indexed by
\(\mathcal P_0\cap Q'\), and \(v\) is indexed by all physical sites
in \(Q'\), modulo constants.  Its dimension is at most
\(3|Q'|\le C_d L_{\mathrm{tile}}^d\le C_d t^{-6d}\).
Its norm and extension map are
\begin{equation}\label{eq:local-pullback}
 \| (\alpha,v)\|_{\rm loc}^2
   =\|\alpha\|_{Q'}^2+b^{-1}\|\nabla v\|_{Q'}^2,
 \qquad
 B_T(\alpha,v)=\chi\big(\alpha+v-\operatorname{avg}_{Q'}v\big).
\end{equation}
The extension is zero outside \(Q'\), and physical values of \(v\) are
repeated on ports.  The subscript \(T\) here identifies the tile.
The symbol \(B_T^*\) denotes the pullback into the dual of the local
space.  A pulled-back symmetric form is measured by the norm of its
representing self-adjoint operator in the local inner product.
The set \(\mathcal P_0\), the cubes, \(\chi\), \(t,b\), the local
norm, and \(B_T\) stay fixed throughout the sweep, including retries.
They may depend on the stage-start history, on which we condition.
Each new correction is extended by zero to \(\mathcal P_0\) at its
time of creation.  The precise record convention below keeps this
reference space fixed at every stopping step.

\begin{lemma}[Local pullback]\label{lem:local-pullback}
The maps in \eqref{eq:local-pullback} satisfy
\[
 W(B_Tf)\le C\|f\|_{\rm loc}^2,
 \qquad
 \|B_T^*y\|_{{\rm loc},*}^2\le CW^*(y)
\]
for neutral charges \(y\).  At a prestate satisfying the inherited
stage-start and current hypotheses of Proposition~\ref{prop:sweep-increments},
the comparison \(D_{\rm cur}(f|_{\rm cur})\le CW(f)\) holds.  Pad each \(L_x\) to the
initial ports.  The pulled-back increments \(B_T^*L_xB_T\) then satisfy the
mean and operator-square-stack bounds of
Proposition~\ref{prop:sweep-increments}, up to a fixed constant, in
the local norm.
\end{lemma}

\begin{proof}
Use the displayed decomposition of \(B_Tf\) in the infimum defining
\(W\).  Its \(\alpha\) part has no larger norm.  The product rule and
cube Poincare inequality give
\[
 \|\nabla[\chi(v-\operatorname{avg}v)]\|^2
 \le C\|\nabla v\|_{Q'}^2
      +CL_{\mathrm{tile}}^{-2}\|v-\operatorname{avg}v\|_{Q'}^2
 \le C\|\nabla v\|_{Q'}^2.
\]
This proves the first estimate, and duality gives the second.

For clarity, let \(\mathcal R\) be the Riesz map from the local space
to its dual, and set \(A_x=\mathcal R^{-1}B_T^*L_xB_T\).  Every
\(L_xy\) is neutral.  Hence
\[
 \sum_x\|A_xf\|_{\rm loc}^2
 \le C\sum_x W^*(L_xB_Tf)
 \le Ct^{\delta/2}D_{\rm cur}((B_Tf)|_{\rm cur})
 \le Ct^{\delta/2}\|f\|_{\rm loc}^2.
\]
Since each \(A_x\) is self-adjoint in this fixed inner product, this
is the scalar identity bound
\[
 \sum_x A_x^2\le C_d t^{\delta/2}\Id_{\mathcal H_T},\qquad
 \left\|\sum_x\EE A_x\right\|_{\mathcal H_T}\le C_d t^\beta.
\]
The mean estimate follows from the corresponding form bound.  Both
assertions hold for every subset of possible next ports; the square
bound holds for arbitrary choices of their hypothetical outcomes.
All maps here act on a fixed initial-port space; deleting a port
merely pads the new matrix by zero there.
\end{proof}

\subsection{Trials and the transfer of local budgets}

Start from a physical state satisfying the complete stage-boundary
bounds of Section~\ref{sec:state}, including self-consistency,
\eqref{eq:comparators}, the broad bounds on \(r,p\), the boundary cap
\(\|K^{\rm phys}\|\le1/3\), the fourth-power row cap \(H_*=t^{1+\delta}\),
and an active physical-site count of at least \(N/2\).  Duplicate it
as in Lemma~\ref{lem:duplication}.  The resulting initial data include
\[
 \widetilde U_0=0,\qquad D_0\asymp W,\qquad
 X_0\le(1+t^\kappa)\diag(p_x^{-2})+C_nt\Lambda^+_\Sigma,
\]
and the full duplicated fourth-power rows are at most \(.75H_*\).
These initial comparisons, the broad bounds on the frozen \(r,p\),
the active physical-site count, and \eqref{eq:local-metric-input} are
standing assumptions for the rest of this section.  Bit-port deletions
leave the physical active and terminal sets unchanged, so the count and
metric remain fixed; the initial comparison data and frozen parameters
also remain fixed.

Process the tiles in a fixed order.  A trial in a tile \(T\) starts
from its current committed numerical pre-tile state and reveals a
uniform permutation of all \(M=L_{\mathrm{tile}}^d\) physical sites in
\(T\), one slot at a time.  A terminal slot is skipped.  At an active
slot the trial deletes the bit port at that site and uses the update of
Lemma~\ref{lem:single-update}.  For each digit stage \(n<m\), tile \(T\),
trial number \(a\ge1\), and physical site \(i\in T\), supply an auxiliary
fair sign \(\varepsilon^{\rm aux}_{n,T,a,i}\).  The auxiliary signs are
independent over these indices and independent of the true digits.
The corresponding physical digit is
\(\sqrt{3s_n/4}\,\varepsilon^{\rm aux}_{n,T,a,i}\).  Also supply
independent uniform permutations \(\pi_{n,T,a}\) for
\(n_0\le n\le m\), each tile \(T\) at that stage, and \(a\ge1\),
independent of both sign families.  Only
the slots and auxiliary signs actually used are revealed.  On a first
digit trial, a coupled active site \(i\) uses
\(\sqrt{3s_n/4}\,\varepsilon_{i,n-n_0+1}\), and a pending site uses
its designated auxiliary digit.  Later digit trials
use their new auxiliary family at every active site.  The actual final
sweep is entered only with no pending site; its first trial at a tile
uses \(\xi_i:=\sum_{k>m-n_0}\lambda2^{-k}\varepsilon_{i,k}\), the whole
remaining true suffix, at every active site and aborts on
failure.  It uses only \(\pi_{m,T,1}\), samples no auxiliary suffix,
and has no retry.

We now specify the numerical records and the test made after a slot.
For each deletion \(x\), retain \(J_x,L_x\) on the fixed initial-port
space \(\mathcal P_0\): pad the new increments by zero at \(x\) and at
ports deleted before that update.  Entries already recorded in an
earlier increment are retained when a later port is deleted.  This is
the convention in the cumulative Schur identities
\eqref{eq:sweep-cumulative-schur}.

There are also three scalar records on fixed coordinate sets.  Let
\(\Delta\mathsf u_i^x\) equal the transformed-lag increment
\(\Delta\widetilde U_i^x\) on the survivors of deletion \(x\), and set
it to zero at \(x\) and at every earlier deleted coordinate.  For
\(\sigma\in\{\mathrm{all},\mathrm{res}\}\), let
\(\Delta\mathsf R_i^{\sigma,x}=R_i^{\sigma,x}/H_*\) be the normalized proxy
increment of Lemma~\ref{lem:sweep-fourth}, with all current columns
tracked when \(\sigma=\mathrm{all}\), and reserve columns tracked when
\(\sigma=\mathrm{res}\).  The latter record is defined only for reserve
rows.  Each proxy sums the fourth-power change over surviving tracked
columns; it omits the favorable loss of the removed column.  Its
increment is zero when its own row is deleted and thereafter.  Thus all
three records have fixed coordinate sets in \(\mathcal P_0\), with a
coordinate frozen at its deletion.  The normalized proxy increments are
the quantities denoted by \(\Delta R_i\) in
\eqref{eq:local-scalar-tails}.

Write \(x_\ell\) for the bit port at slot \(\ell\) when that slot is active, and
interpret each summand below as zero at a terminal slot.  The records
for prefix \(k\) of a trial in \(T\) are
\[
 \begin{aligned}
 J_{T,k}&=\sum_{\ell\le k}J_{x_\ell},&
 L_{T,k}&=\sum_{\ell\le k}L_{x_\ell},\\
 \mathsf u_{T,k}&=\sum_{\ell\le k}\Delta\mathsf u^{x_\ell},&
 \mathsf R^\sigma_{T,k}&=\sum_{\ell\le k}\Delta\mathsf R^{\sigma,x_\ell}.
 \end{aligned}
\]
Every one of these records is zero at the start of each trial, including
a retry.  If an earlier tile \(S\) has been accepted, write
\(J_S,L_S,\mathsf u_S,\mathsf R^\sigma_S\) for the endpoint records of its
accepted trial.  At a current prefix in \(T\), define the accumulated
sweep records by
\begin{equation}\label{eq:local-sweep-records}
 \begin{aligned}
 J^{\rm sw}_{T,k}&=\sum_{S\prec T}J_S+J_{T,k},&
 L^{\rm sw}_{T,k}&=\sum_{S\prec T}L_S+L_{T,k},\\
 \mathsf u^{\rm sw}_{T,k}&=\sum_{S\prec T}\mathsf u_S+\mathsf u_{T,k},&
 \mathsf R^{\sigma,\rm sw}_{T,k}
   &=\sum_{S\prec T}\mathsf R^\sigma_S+\mathsf R^\sigma_{T,k}.
 \end{aligned}
\end{equation}
Only the accepted trial of each earlier tile enters these sums.  A
failed numerical prefix is absent.  In this included numerical record,
a physical site is used as a pivot at most once: the tiles are disjoint
and each trial permutation visits a site once.  The retained information
may contain failed queries at those same physical sites.

Let \(\mathcal P_{\rm cur}\) be the current survivor ports, and let
\(\mathcal P_{\rm res}\) be the reserve ports, which survive the whole
sweep.  Along any included sequence for which the increments are defined,
the lag record has the exact restriction identity
\[
 \mathsf u^{\rm sw}_{T,k}|_{\mathcal P_{\rm cur}}
       =\widetilde U_{\rm cur}.
\]
Indeed, both start at zero and have identical increments on each
survivor; a deleted coordinate of the fixed record keeps its value
immediately before deletion.  The row records have a different role.
For a surviving row, and respectively for a reserve row, they give only
\begin{align*}
 \sum_{j\in\mathcal P_{\rm cur}}|h^{\rm cur}_{ij}|^4
 &\le \sum_{j\in\mathcal P_0}|h^0_{ij}|^4
             +H_*\mathsf R_i^{\mathrm{all},\rm sw},\\
 \sum_{j\in\mathcal P_{\rm res}}|h^{\rm cur}_{ij}|^4
 &\le \sum_{j\in\mathcal P_{\rm res}}|h^0_{ij}|^4
             +H_*\mathsf R_i^{\mathrm{res},\rm sw}.
\end{align*}
Here and below the current \((T,k)\) subscripts are suppressed when
unambiguous.  The inequalities retain the favorable removed-column
loss on the left and omit it from the proxy on the right.

Call the current prefix \emph{valid} when it satisfies both
\eqref{eq:sweep-global-caps} and \eqref{eq:sweep-row-lag-caps}, with
the exact meanings
\begin{equation}\label{eq:local-validity}
 \begin{gathered}
 \|J^{\rm sw}_{T,k}\|\le t^{.10\delta},\qquad
 |L^{\rm sw}_{T,k}|\le t^{.10\delta}D_0,\\
 \max_{i\in\mathcal P_{\rm cur}}
       \sum_{j\in\mathcal P_{\rm cur}}|h^{\rm cur}_{ij}|^4\le H_*,
 \qquad
 \|\widetilde U_{\rm cur}\|_{\ell^\infty(\mathcal P_{\rm cur})}
       \le t^\beta,\qquad \beta=.42\delta.
 \end{gathered}
\end{equation}
The first line tests the accumulated padded corrections on the initial
space, and the second tests the actual current row and transformed lag
on survivors.  The two-sided form inequality for \(L^{\rm sw}_{T,k}\)
is on the initial ports.  The standing data above are not additional
prefix tests.  By the Section~\ref{sec:sweep} Schur and prestate
arguments, validity and those data supply the hypotheses of
Proposition~\ref{prop:sweep-increments}; Lemma~\ref{lem:inverse-row}
supplies its absolute inverse-row bound.

Let \(P_{Q_T}\) be the Euclidean projection onto
\(\mathcal P_0\cap Q_T\), and let \(\mathcal R_T\) be the Riesz map of
\(\mathcal H_T\).  At every trial prefix impose the local budgets
\begin{equation}\label{eq:local-budgets}
 \begin{gathered}
 \|\mathsf u_{T,k}\|_{\ell^\infty(\mathcal P_0\cap Q_T)}
       \le t^{.44\delta},\qquad
 \max_{\sigma\in\{\mathrm{all},\mathrm{res}\}}
       \|\mathsf R^\sigma_{T,k}\|_{\ell^\infty(Q_T)}
       \le t^{.0005\delta},\\
 \|P_{Q_T}J_{T,k}P_{Q_T}\|\le t^{.125\delta},\qquad
 \|\mathcal R_T^{-1}B_T^*L_{T,k}B_T\|_{\mathcal H_T}
       \le t^{.125\delta}.
 \end{gathered}
\end{equation}
For a row record, the norm over \(Q_T\) uses its fixed row indices above
that cube.  The matrix norm is Euclidean and the form norm uses the
fixed local inner product.  These budgets apply to changes accumulated
since this trial began, whereas \eqref{eq:local-validity} uses the full
included sweep record.

After a slot, continue precisely when the local budgets and validity
still hold.  Stop and fail at their first violation, retaining
the update that produces it.  Its preceding prefix satisfies both cap
displays, so Lemma~\ref{lem:local-prospective} supplies the poststate
inverse and a fixed \(I-K'\) invertibility margin.  The transformed lag
and this first possibly failing test are therefore defined without
assuming fine caps at the poststate.
If all \(M\) slots finish without a
violation, accept the trial's numerical changes and endpoint records.
On a failure restore the pre-tile potential values, port set, and
matrices, and discard the failed trial records.  Retain every queried index and answer and
every exposed permutation slot, including the answer producing the
first violation.  A failed first trial marks every active site of the
tile pending, including active sites whose slots were not reached.
Previously terminal sites remain terminal.  When the digit retry
construction is used, its numerical starting state has those restored
values and all active labels in the tile pending.  A later failure
restores that retry starting state and retains its own queried
information; the next retry uses a new auxiliary family and independent
permutation.  At the later cutoff where the construction requires the
pending set to be empty, it begins aborting at any first-trial failure
instead of entering a retry.

Here is the ordinary verification of the input law.  At entry to scale
\(n\), assume the retained-history invariant that each coupled active
site has the first \(n-n_0\) true digits as its partial value and its
remaining suffix is untouched.  A digit trial consumes the next bit;
the final trial at \(n=m\) consumes the whole suffix after those
\(m-n_0\) digits.  The permutation exposes a site before its scalar
answer is queried, and every stopping and acceptance decision is a
function of the resulting retained history.  At a coupled first-trial
slot the required true bit, or final suffix, is therefore unqueried;
at a pending slot the designated auxiliary entry is unused.
Lemma~\ref{lem:conditional-freshness} gives the scalar law conditional
on the retained history before the slot and on the newly exposed slot,
and gives the conditionally uniform order on the unexposed slots of the
new independent permutation.  On success a coupled site consumes that true
input and retains the unused suffix (empty after final completion),
while a pending site stays
pending.  On first-trial failure the numerical rollback leaves every
answer in the history, and the whole-tile marking rule removes all its
active sites from the class at which true freshness is used.  This rule
also conservatively marks unreached active sites, whose own suffixes may
still be untouched.  Retries use only their unused auxiliary entries.
Thus ordinary trials preserve the invariant at every coupled site and
satisfy the conditional interface above after arbitrary previous
failures.  This verification uses the generic reveal lemma, without
conditioning on future acceptance of a trial.

\begin{lemma}[Transfer from tiles to the whole sweep]\label{lem:local-transfer}
Under the standing duplicated stage-start and metric assumptions,
consider the numerical record consisting of the accepted endpoint
records of earlier tiles and one current trial prefix, possibly empty.
Suppose the physical pivots in this collection are distinct and every
included update has a prestate satisfying \eqref{eq:local-validity}.
Assume \eqref{eq:local-budgets} for each accepted endpoint record and
for the current prefix record.  The last included poststate is allowed
to violate a mutable cap.  Then, for sufficiently small \(t\),
\[
 \begin{gathered}
 \|J^{\rm sw}\|\le Ct^{.125\delta}\ll t^{.10\delta},\qquad
 |L^{\rm sw}|\le Ct^{.125\delta}D_0\ll t^{.10\delta}D_0,\\
 \|\mathsf u^{\rm sw}\|_\infty\le Ct^{.44\delta},\qquad
 \max_{\sigma\in\{\mathrm{all},\mathrm{res}\}}
       \|\mathsf R^{\sigma,\rm sw}\|_\infty\le Ct^{.0005\delta}.
 \end{gathered}
\]
Consequently its actual surviving transformed lag and full rows satisfy
\[
 \|\widetilde U_{\rm cur}\|_\infty
       \le Ct^{.44\delta}<t^{.42\delta},\qquad
 \max_{i\in\mathcal P_{\rm cur}}
       \sum_{j\in\mathcal P_{\rm cur}}|h^{\rm cur}_{ij}|^4
       \le(.75+Ct^{.0005\delta})H_*<H_*.
\]
The reserve-row output retained for the endpoint is
\[
 \max_{i\in\mathcal P_{\rm res}}
       \sum_{j\in\mathcal P_{\rm res}}|h^{\rm cur}_{ij}|^4
       \le(\theta^2+Ct^{.0005\delta})H_*.
\]
In particular the current prefix is strictly inside all four mutable
caps in \eqref{eq:local-validity}.
\end{lemma}

\begin{proof}
For each included tile write \(J_T,L_T,\mathsf u_T,\mathsf R^\sigma_T\)
for its accepted endpoint record or its current prefix record.  Every
included update has a valid prestate, so
Lemma~\ref{lem:local-prospective} gives the scalar and matrix tails
through the last candidate update.  Its proof does not require that
candidate poststate to satisfy a fine cap.

For \(J\), bounded overlap of the cubes gives
\[
 \left\|\sum_T1_{Q_T}J_T1_{Q_T}\right\|
 \le C\max_T\|1_{Q_T}J_T1_{Q_T}\|.
\]
Indeed, test against two vectors and apply Cauchy--Schwarz to
\(\sum_T\|1_{Q_T}f\|\|1_{Q_T}g\|\).  No sign assumption on \(J_T\)
is needed.  The discarded entries are negligible by
\eqref{eq:local-matrix-tails} and distinctness of the pivots: an entry
discarded for a tile has an endpoint at distance at least
\(L_{\mathrm{tile}}\) from its pivot, and the remaining exponential
kernel has uniformly bounded row and column sums up to a power of
\(t^{-1}\).  Thus its total operator norm is at most
\(t^{-C}e^{-cuL_{\mathrm{tile}}}\).  The analogous scalar argument uses
bounded overlap at each monitored row and
\eqref{eq:local-scalar-tails} elsewhere.  It gives the stated bounds
for the fixed records \(\mathsf u^{\rm sw}\) and
\(\mathsf R^{\sigma,\rm sw}\), up to the same negligible error.

For the energy form, test an arbitrary split \(f=\alpha+v\) on the
initial ports, and write \(f_T\) for the restrictions of the pair
\((\alpha,v)\) to the local space \(\mathcal H_T\).
On \(Q_T\times Q_T\), differences of \(B_Tf_T\) equal
differences of \(f\).  The error is a tail.  If \(x\) is in the tile
and an endpoint lies outside \(Q_T\), its distance from \(x\) is at
least \(L_{\mathrm{tile}}\).  Extracting part of this distance from
\eqref{eq:local-matrix-tails}, then summing over distinct pivots,
bounds the absolute error edge kernel by
\begin{equation}\label{eq:local-edge-tail}
 t^{-C}e^{-cuL_{\mathrm{tile}}}e^{-cu\dist(i,j)}.
\end{equation}
Convolution sums only change the polynomial prefactor.
For a minimal torus displacement representative \(a\),
\[
 \sum_i|v_i-v_{i+a}|^2\le |a|_1^2\|\nabla v\|^2.
\]
The shell count used in Lemma~\ref{lem:inverse-row} gives, for
\(s=0,2\),
\[
 \sum_a |a|_1^s e^{-cu|a|_1}
 \le C_d\sum_{r\ge0}(1+r)^{d-1+s}e^{-cur}
 \le C_d u^{-d-s},
\]
where for \(s=0\) the zero displacement contributes one.  Thus
the gradient contribution costs at most \(C_du^{-d-2}\), and the
\(\alpha\) contribution at most \(C_du^{-d}\).
It follows that \eqref{eq:local-edge-tail} contributes at most
\[
 t^{-C}e^{-cuL_{\mathrm{tile}}}
       (\|\alpha\|^2+b^{-1}\|\nabla v\|^2).
\]
For the pullback tail of one tile, apply this same kernel estimate to
the zero extensions
\[
 \alpha_T=\chi_T\alpha,\qquad
 v_T=\chi_T(v-\operatorname{avg}_{Q'_T}v).
\]
Its error is bounded by the exponential prefactor times
\(\|\alpha_T\|^2+b^{-1}\|\nabla v_T\|^2\), which is at most
\(C\| (\alpha,v)\|_{{\rm loc},T}^2\) by the product and Poincare
estimate in Lemma~\ref{lem:local-pullback}.  Summing these local
energies costs only bounded overlap.  Since
\(uL_{\mathrm{tile}}\asymp t^{-3}\), all these errors are smaller than
any power of \(t\).  The local form budgets give
\[
 \sum_T |(B_T^*L_TB_T)(f_T)|
 \le t^{.125\delta}\sum_T\|f_T\|_{\rm loc}^2
 \le Ct^{.125\delta}
       (\|\alpha\|^2+b^{-1}\|\nabla v\|^2).
\]
Combining this with the two tail estimates and taking the infimum over
splits yields
\[
 |L^{\rm sw}|\le Ct^{.125\delta}W
       \le Ct^{.125\delta}D_0\ll t^{.10\delta}D_0.
\]
Here the last comparison is the initial comparison
\(D_0\asymp W\), not an assumed post-update comparison.
The preceding overlap bounds similarly give
\(\|J^{\rm sw}\|\le Ct^{.125\delta}\ll t^{.10\delta}\) and
\(\|\mathsf u^{\rm sw}\|_\infty\le Ct^{.44\delta}\ll t^{.42\delta}\).
Restricting the latter record to survivors gives the actual transformed
lag bound.  The total normalized row-proxy increase is at most
\(Ct^{.0005\delta}\).  Its upper-bound relation with actual row growth,
together with the initial full-row bound \(.75H_*\) and the initial
reserve-row bound \(\theta^2H_*\), proves the two row conclusions.
All the conclusions therefore hold through the candidate update.
\end{proof}

At a first possible failure with no local crossing, this transfer will
force the four mutable caps to remain strict.  We now estimate the
probability of a local crossing in one tile.

\subsection{A maximal concentration estimate}

We will repeatedly apply the following finite-dimensional estimate
inside a neighborhood whose dimension is polynomial in $t^{-1}$ for
fixed $d$. The variance envelopes below are predictable scalars
dominating the conditional squares in matrix order. This precise
version follows from the finite-dimensional Golden--Thompson
inequality, as the proof shows. For the general matrix-martingale
framework, compare \cite[Theorem~1.2 and Corollary~1.3]{Tropp2011};
the scalar-envelope version used here is proved below.

\begin{lemma}\label{lem:matrix-concentration}
Let $Y_k$ be self-adjoint $D_s$-dimensional matrix martingale differences
adapted to $(\mathcal F_k)$, with $\|Y_k\|\le J$. Suppose
\[
 \EE(Y_k^2\mid\mathcal F_{k-1})\le v_k\Id,
 \qquad \sum_k v_k\le V
\]
pathwise, where $v_k$ is nonnegative and $\mathcal F_{k-1}$-measurable,
and $J,V$ are deterministic. Then for $a>0$, with an absolute $c>0$,
\[
 \PP\left(\sup_n\left\|\sum_{k\le n}Y_k\right\|>a\right)
 \le2D_s\exp\left[-c\min\left(\frac{a^2}{V},\frac aJ\right)\right].
\]
The zero-variance case is understood by continuity. A stopped process
may be continued by zero increments.
\end{lemma}

\begin{proof}
For $u>0$ with $uJ\le1$, functional calculus gives
$e^{uY_k}\le\Id+uY_k+Cu^2Y_k^2$. Taking conditional expectations,
using the martingale-difference property, and applying
$1+x\le e^x$ yields
$\EE(e^{uY_k}\mid\mathcal F_{k-1})\le e^{Cu^2v_k}\Id$.
The Golden--Thompson inequality~\cite{Golden1965,Thompson1965},
in the finite-Hermitian form stated in
\cite[Equation~(1.1)]{ForresterThompson2014}, says
$\Tr e^{B+C}\le\Tr(e^Be^C)$ for
finite-dimensional Hermitian $B,C$. Applied conditionally, and using
that $v_k$ is a predictable scalar, it shows that
\[
 Z_n=\Tr\exp\left(u\sum_{k\le n}Y_k-Cu^2\sum_{k\le n}v_k\Id\right)
\]
is a nonnegative supermartingale with $Z_0=D_s$.
At the first upper eigenvalue crossing of $a$, it is at least
$\exp(ua-Cu^2V)$. Stopping at that crossing and a finite horizon bounds
its probability by $D_s\exp(-ua+Cu^2V)$. Let the horizon increase, choose
$u$ to be a small constant times $\min(a/V,1/J)$, and repeat with $-Y_k$.
The union of the two crossings gives the result.
\end{proof}

A rectangular or complex increment $B$ is handled by the Hermitian
matrix $\left(\begin{smallmatrix}0&B\\B^*&0\end{smallmatrix}\right)$.
Its square requires bounds for both $BB^*$ and $B^*B$.
Conditional centering doubles the jump bound at most. It does not
increase either conditional square: if $\bar B=\EE(B\mid\mathcal F)$,
then
\[
 \EE[(B-\bar B)(B-\bar B)^*\mid\mathcal F]
 =\EE[BB^*\mid\mathcal F]-\bar B\bar B^*
 \le\EE[BB^*\mid\mathcal F],
\]
and the same calculation holds in the other order. For vector
increments, the squared Euclidean norm bounds both dilation blocks.
All spaces and their inner products are fixed at the start of each
trial; removed indices are padded by zeros.

Here is how scalar envelopes arise in a randomly ordered trial. If
there are $r$ equally likely next slots and their possible square-matrix
increments $B_x$ satisfy, for every vector $v$ and every collection of outcomes,
\[
 \sum_x\bigl(\|B_xv\|^2+\|B_x^*v\|^2\bigr)
 \le V_*\|v\|^2,
\]
then averaging first over the outcome in each slot and then over its
uniform selection bounds each conditional dilation square by
$(V_*/r)\Id$. Only the conditional marginal laws of these hypothetical
outcomes are needed. Over $M$ slots their envelopes sum to at most
$V_*\sum_{r=1}^M r^{-1}$. For $M\le t^{-C_d}$ this costs
$O_d(\log(1/t))$. The localized energy operators and vector arrays used
later have the corresponding square-stack bounds.
In every application below, the
predictable variance sum is bounded through the first prospective
budget-breaking update. That increment is included whenever the trial's
tests permit another update at the preceding prefix; this decision is
predictable, and all subsequent increments are zero. Thus the concentration
hypotheses are verified
before the outcome that might violate a budget is observed.

\subsection{One trial and its retries}

\begin{lemma}[One trial]\label{lem:local-trial}
Under the standing duplicated stage-start and metric assumptions,
suppose all earlier tiles were accepted under the trial rule and the
present trial starts with their accepted matrix state and accumulated
records satisfying \eqref{eq:local-validity}.  Its active labels may reflect earlier failed attempts at the
current tile.  Assume the conditional slot and scalar
interface stated above.  Uniformly in the full retained history at
this trial start, the conditional probability that the next trial fails
is at most \(\exp\{-t^{-c}\}\), for a fixed \(c>0\).
\end{lemma}

\begin{proof}
Condition on the entire retained history at the trial start.  There are
\(M=L_{\mathrm{tile}}^d\asymp_d t^{-6d}\) permutation slots. While the
trial is running, let \(\mathcal F_{k-1}\) be the full retained history
immediately before slot \(k\). It includes all earlier queries, their answers,
exposed permutation slots and decisions, but not the next slot or its
answer. Given \(\mathcal F_{k-1}\), the next member is uniform among
the \(r_k=M-k+1\) unexposed physical sites.  A terminal slot has zero
increment.  At each other possible slot, the scalar law conditional on
\(\mathcal F_{k-1}\) and that newly exposed slot satisfies
\eqref{eq:sweep-digit-moments}.  The current numerical
prestate is \(\mathcal F_{k-1}\)-measurable.

Let \(\tau\) be the first prefix violating a local budget or one of the
four mutable caps, and put \(\tau=M+1\) if there is none.  The multiplier
\(s_k=\mathbf1_{\{k\le\tau\}}\) is predictable: it asks only whether
prefixes through \(k-1\) satisfied the local budgets and validity.
When \(s_k=1\), evaluate the next raw increment at that valid prestate;
when \(s_k=0\), set the stopped increment to zero without another query.
This defines the stopped increment as \(s_k\) times the raw increment
where it is evaluated.  Thus the first breaking update is included,
and every retained update has a valid prestate. After stopping, the formal
remaining increments are zero and \(\mathcal F_k\) stays equal to the history
at stopping. The actual retained filtration continues through rollback,
marking, and later trials. The poststate of the
first breaking update need not satisfy the fine caps;
Lemma~\ref{lem:local-prospective} supplies all decay estimates needed
through the breaking update.  The next-slot law and its count \(r_k\)
are used only when \(s_k=1\); after stopping every variance envelope is zero.

We verify the scalar identity variance hypothesis of
Lemma~\ref{lem:matrix-concentration} separately for every monitored
quantity.  Put \(\ell=1+\log(1/t)\), let \(\EE_{k-1}\) denote
conditional expectation given \(\mathcal F_{k-1}\), and use the
Hermitian dilation
\[
 \mathsf H(B)=\begin{pmatrix}0&B\\B^*&0\end{pmatrix},\qquad
 \mathsf H(B)^2=\diag(BB^*,B^*B).
\]
All operators below act on the fixed spaces and scalar coordinate sets
defined above.

For a monitored lag coordinate, take
\(\mathsf H(\Delta\mathsf u_i^x)\) on \(\mathbb C^2\).
Its square is \(|\Delta\mathsf u_i^x|^2\Id_2\), so the lag
stack of Proposition~\ref{prop:sweep-increments} gives
\(C_d\ell^2t^\delta\Id_2\) after summing over all remaining
possible bit ports.  The record has zero increment at its own
deletion and thereafter.

For \(J\), write \(P_Q\) for the fixed Euclidean projection onto
\(\mathcal P_0\cap Q\) and \(C_x=P_QJ_xP_Q\) on that space.
Both square orders are controlled:
\[
 \sum_x C_xC_x^*\le P_Q\Big(\sum_xJ_xJ_x^*\Big)P_Q
       \le C_dt^{\delta/2}\Id_Q,\qquad
 \sum_x C_x^*C_x\le C_dt^{\delta/2}\Id_Q.
\]
Consequently the sum of squares of \(\mathsf H(C_x)\) is bounded
by \(C_dt^{\delta/2}\) times the identity on the doubled local
port space.  Zero padding makes these compressions meaningful even
after ports are deleted.

For \(L\), use the self-adjoint Riesz operator
\(A_x=\mathcal R_T^{-1}B_T^*L_xB_T\) on \(\mathcal H_T\).
Lemma~\ref{lem:local-pullback} gives directly
\(\sum_x A_x^2\le C_dt^{\delta/2}\Id_{\mathcal H_T}\).
For either normalized row proxy, its increment \(\Delta\mathsf R_i^{\sigma,x}\)
is real; Proposition~\ref{prop:sweep-increments} gives
\(\sum_x|\Delta\mathsf R_i^{\sigma,x}|^2\le C_dt^\kappa\).  This is an identity
bound on its fixed one-dimensional real space, again with zero
increments from the row's deletion onward.

Each of these pathwise stacks holds over any remaining subset and
for every collection of allowed hypothetical scalar outcomes.
They can therefore be averaged under any joint law with the required
conditional marginals.  Independence among hypothetical outcomes at
different possible next ports is unnecessary.  Averaging over the
uniform next slot divides the relevant stack by \(r_k\).
For each of the four self-adjoint stopped raw increments \(Z_k^a\),
center the joint slot-and-scalar outcome by setting
\(Y_k^a=Z_k^a-\EE_{k-1}Z_k^a\).  Then
\[
 \EE_{k-1}(Y_k^a)^2
 =\EE_{k-1}(Z_k^a)^2-(\EE_{k-1}Z_k^a)^2
 \le v_k^a\Id_a,
\]
with the explicit predictable scalar choices
\[
 \begin{aligned}
 v_k^{\rm lag}&=C_ds_k\ell^2t^\delta/r_k,&
 v_k^J&=C_ds_kt^{\delta/2}/r_k,\\
 v_k^L&=C_ds_kt^{\delta/2}/r_k,&
 v_k^{\rm row}&=C_ds_kt^\kappa/r_k.
 \end{aligned}
\]
Here \(\Id_a\) is the identity on the indicated fixed space.
Centering increases the jump bound by at most a factor two.  The
harmonic sum obeys
\[
 \sum_{k=1}^M r_k^{-1}\le1+\log M\le C_d\ell.
\]
Thus the deterministic total variance envelopes, centered jump
bounds, and budgets, up to constants depending only on \(d\), are
\[
\begin{array}{c|ccc}
 \text{quantity}&V&\text{jump}&\text{budget}\\\hline
 \widetilde U_i&\ell^3t^\delta&\ell t^{\delta/2}&t^{.44\delta}\\
 J,\ B_T^*LB_T&\ell t^{\delta/2}&t^{\delta/4}&t^{.125\delta}\\
 \text{normalized row proxy}&\ell t^\kappa&t^{\kappa/2}&t^{.0005\delta}
\end{array}
\]
The same averaging of the arbitrary-subset mean estimates bounds the
sum of norms of conditional means by
\[
 \begin{aligned}
 B_{\rm lag}&=C_d[\ell^2(t^{2\beta}+\sqrt t)
                         +\ell^3t^{\delta/2}],\\
 B_J+B_L&\le C_d\ell t^\beta,\qquad
 B_{\rm row}\le C_d\ell t^\kappa.
 \end{aligned}
\]
These are respectively \(o(t^{.44\delta})\), \(o(t^{.125\delta})\),
and \(o(t^{.0005\delta})\), so each is smaller than half its budget
for sufficiently small \(t\).
Apply Lemma~\ref{lem:matrix-concentration} to the centered sums at
half-budget.  The exponent margins for squared budget divided by \(V\) and for
budget divided by the jump bound are, respectively, \(.25\delta\) and
\(.125\delta\) for \(J,L\), \(.12\delta\) and \(.06\delta\) for
the transformed lag, and \(.004\delta\) and \(.002\delta\) for a
normalized row proxy, using \(\kappa=.005\delta\).  For the lag, the
two ratios are bounded below by
\[
 ct^{-.12\delta}/\ell^3,\qquad ct^{-.06\delta}/\ell.
\]
All these ratios tend to infinity by a positive power of \(t^{-1}\),
apart from logarithms.

The \(J\) dilation has dimension \(2|\mathcal P_0\cap Q|\le
C_dt^{-6d}\), and \(\dim\mathcal H_T\le C_dt^{-6d}\).  Each lag
test has dimension two and each row-proxy test dimension one; their
total number is at most \(C_dt^{-6d}\).  The sum of all dimension
prefactors in the concentration bound is consequently at most
\(C_dt^{-6d}\), independent of \(N\).  The union over these local
tests has probability at most \(e^{-t^{-c}}\), after decreasing
\(c>0\) and increasing the dimension-dependent starting scale.
There is no union over the entire torus in this single-trial bound.

Suppose now that the trial fails at prefix \(\tau\), but no local
budget is violated there.  Earlier accepted trials never crossed either
test, and all earlier current prefixes satisfy validity.  Thus
every included update, including the one at \(\tau\), has a valid
prestate.  The accepted earlier endpoint records satisfy the
local budgets, and by hypothesis so does the current record at
\(\tau\).  Their physical pivots are distinct by the record convention
following \eqref{eq:local-sweep-records}.  Lemma~\ref{lem:local-transfer}
therefore puts the poststate at \(\tau\) strictly inside all four
mutable caps, contradicting the definition of \(\tau\).  Every trial
failure thus includes a local budget crossing; a simultaneous violation
of a local budget and a mutable cap is already counted in that event.  The local
crossing estimate proves the lemma.
\end{proof}

The same estimate applies after any number of failed retries whenever
each retry is supplied with a new independent permutation and fresh
scalar inputs satisfying the conditional moment interface.  The indexed
auxiliary signs supply these inputs for the digit sweeps.  Each retry
begins from the restored valid numerical pre-tile state, while its law
is taken with respect to the larger retained history.  Let
\(\mathscr F_T\) be the retained pre-tile history and set
\(p_{\rm fail}=\exp\{-t^{-c}\}<1\). For each fixed attempt number \(a\),
stop at the first of the start of attempt \(a\) and completion of the tile
by an earlier attempt. This stopping time of the retained filtration is almost surely
finite: before it there are at most \(a-1\) attempted trials, each with at
most \(M\) slots. Assign the \(a\)-th failure indicator value zero when an
earlier attempt has completed the tile. The conditional failure bound then
holds at every such stopped history, using Lemma~\ref{lem:local-trial} when
attempt \(a\) is scheduled and the zero indicator otherwise. Successive
conditioning gives
\[
 \PP(\text{the next }k\text{ trials all fail}\mid\mathscr F_T)
 \le p_{\rm fail}^k=e^{-kt^{-c}}.
\]
Hence a tile has infinitely many failures with probability zero.  There
are finitely many tiles at a fixed sweep, so any such supplied retry
process terminates almost surely.  The actual construction uses this
retry conclusion for digit sweeps before the cutoff and uses the
first-trial estimate after the cutoff, including the final sweep.
This argument conditions only
on histories already reached; no estimate within a trial conditions on
its eventual acceptance.

\subsection{The endpoint of an ordinary sweep}

We have proved that accepted locally monitored trials maintain all four
mutable caps.  The remaining task is to express a completed sweep in the physical
normalization at the smaller damping scale.

\begin{proposition}[Ordinary sweep]\label{prop:ordinary-sweep}
Suppose a committed stage state satisfies the complete stage-boundary
bounds of Section~\ref{sec:state}, including the broad bounds and
comparators \eqref{eq:comparators}, the self-consistency relation
\(r_i=g_{ii}\), the fourth-power bound \(H_*=t^{1+\delta}\), the boundary
cap \(\|K_{\rm old}\|\le1/3\), and
\eqref{eq:local-metric-input}.  Assume also that at least \(N/2\)
physical sites are active, where \(N\) is the torus volume, and assume
the conditional slot and scalar interface above.
Let \(\theta\in[.25,.625]\) be the reserve fraction, so the new
damping scale is \(t'=\theta t\).  This includes
\(\theta=t_{n+1}/t_n\) at ordinary stages and \(\theta=1/2\) in the
final residual sweep.  Set \(w=1-\theta\) and use port weights
\(k_{\rm bit}=w^{1/4}\), \(k_{\rm reserve}=\theta^{1/4}\) in
Lemma~\ref{lem:duplication}.  Queried variables in port normalization
are real variables satisfying the conditional moment and size
assumptions above.  If every failed trial is followed by a retry supplied
with a new independent permutation and fresh inputs satisfying this
interface, the trial process terminates almost surely.  The indexed
auxiliary signs supply these retries for digit sweeps; the actual final
sweep instead aborts at its first failure.  Every completed sweep leaves
the terminal set unchanged and has the following endpoint bounds, with \(r,p\) frozen
in physical normalization:
\begin{align*}
 X_{\rm new}
  &=(1+O(t^{.10\delta}))
       \big[\theta X_{\rm old}
               +(1-\theta)\diag(p^{-2})\big],\\
 \|K_{\rm new}\|&\le\|K_{\rm old}\|+t^{.10\delta},\\
 \sup_i|g^{\rm new}_{ii}-r_i|&\le Ct^\beta,\\
 \sup_i\sum_j|h^{\rm new}_{ij}|^4
   &\le\big(\theta^2+Ct^{.0005\delta}\big)t^{1+\delta}
     <.7(\theta t)^{1+\delta}.
\end{align*}
The first line denotes two-sided quadratic-form bounds on neutral
charges.  A first trial at any tile fails conditionally with probability
at most \(e^{-t^{-c}}\).

In particular, with \(b'=\theta b\), the lower and low-frequency
constants in \eqref{eq:comparators} change by relative
\(1+O(t^{.10\delta})\); the rough diagonal upper constant becomes at
most \((1+O(t^{.10\delta}))(\theta C_0+C)\); and the sharp diagonal
upper coefficient becomes
\(1+\theta t^\kappa+O(t^{.10\delta})\).
\end{proposition}

\begin{proof}
Lemma~\ref{lem:local-trial} and successive conditioning prove the
stated conditional retry termination.  Lemma~\ref{lem:local-transfer}
and the stopping rule maintain the mutable caps along every accepted
sweep, using Lemma~\ref{lem:local-prospective} through each candidate
breaking update.  The exact padded Schur
identities then give the \(K\) estimate.  The two-sided \(L\) cap keeps
\(D_0+L^{\rm sw}\) comparable to \(D_0\), with only the constant kernel,
so Schur duality gives multiplicative transport of \(X\) on
survivor-supported neutral charges.  The only surviving ports
are reserve ports, and on those ports
\[
 h_{\rm cur}=\sqrt{\theta t}\,g^{\rm phys,new},\qquad
 p_{\rm cur}=\sqrt\theta\,p^{\rm phys},\qquad
 X^{\rm phys,new}=\theta X_{\rm cur}.
\]
For a neutral physical charge \(y\) placed on the reserve ports, the
exact initial identity \eqref{eq:sweep-reserve-target} reads
\[
 \theta X_0(y_{\rm res})=\theta X_{\rm old}(y)
             +(1-\theta)\sum_i\frac{|y_i|^2}{p_i^2}.
\]
Combining these two formulas with the multiplicative transport proves
the first assertion.  Applying each old comparison separately yields
the three stated new coefficients.  In particular the sharp
coefficient follows from
\(\theta(1+t^\kappa)+(1-\theta)=1+\theta t^\kappa\).

Lemma~\ref{lem:sweep-prestate} gives the actual virtual diagonal
difference \(g^{\rm cur}_{ii}-r_i=O(t^\beta)\); division by the reserve
factor \(\sqrt\theta\) preserves that bound physically.  The reserve-row
conclusion of Lemma~\ref{lem:local-transfer} is
\((\theta^2+Ct^{.0005\delta})H_*\): its fixed proxy bounds the possible
increase after discarding favorable removed-column losses.  At the endpoint the virtual \(h\) is
already the new physical \(h\).  Finally,
\(\sup_{.25\le\theta\le.625}\theta^{1-\delta}<.7\), so the strict
last inequality holds for small \(t\).
\end{proof}

After the cutoff the actual construction requires that no active site
is pending and aborts at any later first-trial failure.  On a retained
history that reaches and completes the final residual sweep without
abort, every accepted final update therefore uses the whole true suffix;
no auxiliary retry contributes to the certificate endpoint.  Its law is
used at its retained pre-query history, before final acceptance is known.
With \(t=2\eta\) and \(w=1/2\), the physical update is
\(\xi_i+\eta r_i\): it inserts the complete true suffix and changes the
artificial term from \(-2\eta r_i\) to \(-\eta r_i\).  Terminal sites
already have their true values.  Thus all physical potentials at this
endpoint are the original true values, while the remaining active sites
are still active and retain the damping \(-\eta\diag_A r\).
